% ============================================================================
% Fichier  : partie6/P06_C20_mesurer.tex : FORN (LPI342)
% Sources : ISO/IEC 27043 ; cadre OUI/NON/NEK (Minka). Opération TSANAGA.
% ============================================================================

\chapter{Mesurer et Progresser}
\label{chap:mesurer}
\chaptermark{Mesurer et conclure}

\epigraph{%
    « La qualité d'une enquête se mesure aussi\\
    à l'honnêteté de ses incertitudes. »%
}{Principe du cours}

\niveauChapitre{Avancé}{L'ensemble du cours~; notions d'auditologie.}

\begin{boxobjectifs}
\begin{itemize}
    \item Définir la qualité d'une enquête forensique.
    \item Choisir des indicateurs honnêtes.
    \item Appliquer le verdict OUI, NON, NEK.
    \item Conduire une revue par les pairs.
    \item Ouvrir sur la recherche.
\end{itemize}
\end{boxobjectifs}

\bigskip

% ============================================================================
\section*{Question directrice : l'enquête TSANAGA était-elle « bonne » ?}
% ============================================================================

Une enquête se juge non à sa conclusion, mais à sa méthode, sa traçabilité et
l'honnêteté de ses limites. Ce chapitre clôt le cours sur la mesure de la
qualité.

\medskip
\textbf{Plan du chapitre.}
\begin{enumerate}
    \item La qualité d'une enquête (section~\ref{sec:c20-qualite}).
    \item Des indicateurs honnêtes (section~\ref{sec:c20-indicateurs}).
    \item Le verdict OUI / NON / NEK (section~\ref{sec:c20-verdict}).
    \item Revue par les pairs et recherche (section~\ref{sec:c20-pairs}).
\end{enumerate}

\separateur

% ============================================================================
\section{La Qualité d'une Enquête}
\label{sec:c20-qualite}
% ============================================================================

\begin{boxdef}
\termecle{Qualité d'une enquête forensique}~: aptitude de l'enquête à produire
des conclusions \emph{fiables} et \emph{opposables}, par une méthode tracée,
reproductible, et une honnêteté assumée sur les limites. Une « bonne »
conclusion issue d'une mauvaise méthode ne vaut rien.
\end{boxdef}

% ============================================================================
\section{Des Indicateurs Honnêtes}
\label{sec:c20-indicateurs}
% ============================================================================

\begin{tableaucours}{L{4.8cm}L{7.4cm}}{%
    \tableauHeader{Indicateur} &
    \tableauHeader{Ce qu'il dit}
}
Traçabilité des actes & Chaque geste documenté et reproductible \\
Intégrité des pièces & Condensats vérifiés à chaque étape \\
Part des conclusions corroborées & Faits tenus par plusieurs sources \\
Clarté des limites & NEK déclarés, incertitudes assumées \\
Tenue en contradictoire & Pièces résistant à la contestation \\
\end{tableaucours}
\vspace{-0.8em}
\begin{center}{\small\itshape Tableau~20.1. Indicateurs de qualité d'une enquête}\end{center}

\begin{boxalert}
\textbf{Indicateur trompeur.} « L'enquête a conclu à la culpabilité. »

\smallskip
\textbf{Règle~:} la conclusion n'est pas un indicateur de qualité~; l'expert
ne conclut d'ailleurs pas à la culpabilité. Ce qui se mesure, c'est la méthode
et l'honnêteté, non le résultat judiciaire.
\end{boxalert}

% ============================================================================
\section{Le Verdict OUI / NON / NEK}
\label{sec:c20-verdict}
% ============================================================================

\begin{boxdef}
\termecle{Verdict à trois valeurs}~: pour chaque question d'enquête, la réponse
est \sigle{OUI} (établi, avec preuve), \sigle{NON} (réfuté, avec preuve) ou
\sigle{NEK} (indéterminé, faute de preuves). Le NEK n'affaiblit pas l'enquête~:
il en garantit l'honnêteté.
\end{boxdef}

\begin{tableaucours}{L{2.0cm}L{4.8cm}L{5.4cm}}{%
    \tableauHeader{Verdict} &
    \tableauHeader{Question (TSANAGA)} &
    \tableauHeader{Fondement}
}
OUI & Le fichier a-t-il été copié vers un amovible~? & Disque et artefacts \\
OUI & A-t-il été téléversé vers un service~? & Réseau et synchronisation \\
NEK & Que contient le second conteneur chiffré~? & Clé introuvable \\
\end{tableaucours}
\vspace{-0.8em}
\begin{center}{\small\itshape Tableau~20.2. Le verdict OUI / NON / NEK appliqué~\cite{minka2026auditologie}}\end{center}

% ============================================================================
\section{Revue par les Pairs et Recherche}
\label{sec:c20-pairs}
% ============================================================================

\begin{boxscenario}{Opération TSANAGA. La relecture}
\textbf{Constatation.} Avant remise, le rapport est relu par un second expert
du cabinet, qui vérifie méthode, condensats, et la juste attribution des
certitudes.

\medskip
\textbf{Analyse} (confiance élevée). La revue par les pairs détecte les
conclusions trop assurées et les NEK oubliés. Elle est, pour l'enquête
forensique, l'équivalent du contrôle qualité~: elle améliore la fiabilité et
l'opposabilité.
\end{boxscenario}

La recherche prolonge le cours~: automatisation de l'analyse, forensique du
nuage et du mobile en contexte africain, et application du cadre de
l'auditologie à la réduction de l'incertitude dans l'enquête numérique.

\begin{boxmetier}{Expert judiciaire et analyste forensique}
\textbf{Se faire relire.} L'expert ne remet jamais un rapport sans relecture
par un pair~: un regard neuf voit les suppositions déguisées en faits. La
qualité se construit à deux.
\end{boxmetier}

\separateur

\begin{boxresume}
\textbf{Ce qu'il faut retenir de ce chapitre~:}
\begin{itemize}
    \item \textbf{La qualité} tient à la méthode et à l'honnêteté, non à la
          conclusion.
    \item \textbf{Des indicateurs}~: traçabilité, intégrité, corroboration,
          limites.
    \item \textbf{OUI, NON, NEK}~: le verdict honnête.
    \item \textbf{La revue par les pairs} est le contrôle qualité de
          l'enquête.
    \item \textbf{La recherche} prolonge la discipline.
\end{itemize}
\end{boxresume}

\begin{boxinfo}
\textbf{Points de contrôle de l'élève.} Je sais~: définir la qualité d'une
enquête~; choisir des indicateurs honnêtes~; appliquer OUI/NON/NEK~; conduire
une revue par les pairs~; situer la recherche.
\end{boxinfo}

\bigskip

\begin{boxexo}[title={\ding{45}\quad Exercice~20.1. Qualité \niveaubase}]
À quoi tient la qualité d'une enquête forensique~?
\end{boxexo}

\begin{boxexo}[title={\ding{45}\quad Exercice~20.2. Trois valeurs \niveaubase}]
Donnez le sens de OUI, NON et NEK.
\end{boxexo}

\begin{boxexo}[title={\ding{45}\quad Exercice~20.3. Trompeur \niveaubase}]
Pourquoi « l'enquête a conclu à la culpabilité » n'est-il pas un indicateur de
qualité~?
\end{boxexo}

\begin{boxexo}[title={\ding{45}\quad Exercice~20.4. Classer \niveaumoyen}]
Attribuez OUI, NON ou NEK aux trois questions de TSANAGA du tableau.
\end{boxexo}

\begin{boxexo}[title={\ding{45}\quad Exercice~20.5. Pairs \niveaumoyen}]
Que détecte une revue par les pairs~?
\end{boxexo}

\begin{boxexo}[title={\ding{45}\quad Exercice~20.6. Méthode ou résultat \niveaumoyen}]
Pourquoi mesure-t-on la méthode et non le résultat judiciaire~?
\end{boxexo}

\begin{boxexo}[title={\ding{45}\quad Exercice~20.7. Indicateurs \niveaumoyen}]
Choisissez trois indicateurs de qualité pour le cabinet NYONG SÉCURITÉ.
\end{boxexo}

\begin{boxexo}[title={\ding{45}\quad Exercice~20.8. NEK \niveaumoyen}]
Pourquoi le NEK garantit-il l'honnêteté de l'enquête~?
\end{boxexo}

\begin{boxexo}[title={\ding{45}\quad Exercice~20.9. Grille de qualité \niveauexpert}]
Construisez une grille d'évaluation de la qualité d'une enquête forensique.
\end{boxexo}

\begin{boxexo}[title={\ding{45}\quad Exercice~20.10. Question de recherche \niveauexpert}]
Formulez une question de recherche sur la forensique du mobile en contexte
africain.
\end{boxexo}

\bigskip

\begin{boxinfo}
\textbf{Pour aller plus loin.}
\begin{itemize}
    \item ISO/IEC 27043~\cite{iso27043}~; cadre de
          l'auditologie~\cite{minka2026auditologie}.
\end{itemize}
\end{boxinfo}

\begin{boxinfo}
\textbf{Fin du cours.} De la scène saisie au verdict mesuré, FORN aura montré
qu'une enquête numérique établit \emph{et} prouve~: acquérir sans altérer,
analyser chaque domaine, reconstituer un faisceau, et restituer avec rigueur,
sans jamais confondre ce que l'on constate et ce que l'on en déduit.
\end{boxinfo}
